\chapter{Полный стационарный фон тензорного вихря}
\label{ch:v6-bosonic-defect-full-tensor-stationary-background-gate}

\section{Пропущенное уравнение}

Скрученная область определения полного гессиана уже выведена, но прежний
фон содержал только три профиля $(K,a,b)$ и фиксировал $Q=Q_*$. Проверка
первой вариации показала, что такой фон не является стационарной точкой
полного действия: остаток уравнения $Q$ в сердцевине ненулевой и направлен
вдоль самого вакуумного тензора $Q_*$. Поэтому до спектрального расчёта
необходимо расширить фон, а не интерпретировать малые уровни гессиана.

Требование сначала решить полные фоновые уравнения, а затем отделять
физические и калибровочные моды, согласуется со стандартным анализом
вихревых решений и их нулевых мод
\cite{stationary-background-alonso,stationary-background-a4-vortex}.

\section{Минимальный четырёхпрофильный анзац}

Остаток, ортогональный $Q_*$, численно меньше $3\cdot10^{-10}$. Следовательно,
симметрия требует лишь одной дополнительной функции:
\begin{equation}
 Q(r)=q(r)Q_*,
 \qquad
 T(r)=b(r)T_0+a(r)T_3,
 \qquad
 B_\theta=-\frac{K(r)}3J_0.
\end{equation}
Новые свободные коэффициенты не вводятся. Единый условный след даёт
\begin{align}
 A&=\frac{128}{243},&
 B&=\frac{160}{243},&
 D_Q&=\frac13\|Q_*\|^2=0.1675239755\ldots,&
 G&=\frac2{27}.
\end{align}

Совместная краевая задача имеет условия
\begin{align}
 K(0)&=a(0)=0,& b'(0)&=q'(0)=0,\\
 K(\infty)&=a(\infty)=b(\infty)=q(\infty)=1.
\end{align}
Потенциал не подгоняется: это точное ограничение прежнего полного
потенциала $V(Q,T)$ на четырёхпрофильный анзац.

\section{Численное решение}

Метод краевой задачи сошёлся на $2099$ адаптивных узлах. Максимальный
относительный остаток равен
\begin{equation}
 1.9851\cdot10^{-7},
\end{equation}
а граничный остаток равен $2.1\cdot10^{-27}$. В сердцевине получено
\begin{align}
 b(0)&=1.0913000372\ldots,\\
 q(0)&=0.9666849845\ldots,\\
 a'(0)&=2.9489441346\ldots.
\end{align}
Таким образом, квадрупольный порядок не исчезает в ядре, а ослабляется на
$3.3315\%$. Этого небольшого изменения достаточно, чтобы восстановить
пропущенное уравнение Эйлера--Лагранжа.

Полное безразмерное натяжение равно
\begin{equation}
 \mathcal T_{4}=1.5744526907\ldots.
\end{equation}
Для трёхпрофильного исправленного фона было
$\mathcal T_{3}=1.5772484656\ldots$, поэтому релаксация $Q$ понижает
энергию на
\begin{equation}
 \mathcal T_3-\mathcal T_4=0.0027957750\ldots.
\end{equation}
Вклад радиального профиля $q$ мал,
\begin{equation}
 E_{r,Q}=0.0002529346\ldots,
\end{equation}
но его вариационное значение принципиально. Вириальный остаток равен
$6.27\cdot10^{-9}$.

\begin{theorem}[Стационарный фон полного анзаца]
Без введения нового параметра существует конечный четырёхпрофильный вихрь
$(K,a,b,q)$, который решает все фоновые уравнения в совместном
$Z_3$-инвариантном секторе. Замороженное приближение $Q=Q_*$ заменяется
единственным радиальным профилем $q(r)$.
\end{theorem}

\section{Нулевая гипотеза и стоп-критерий}

Теперь устранён барьер ненулевой первой вариации. Нулевая гипотеза следующего
гейта: на новом стационарном фоне один из двух нетривиальных скрученных
секторов всё же содержит сходящуюся отрицательную моду.

Стоп-критерий: отрицательный предел после сгущения полярной сетки закрывает
прямую нить как седло полного действия. Если же остаются только две
переносные нулевые моды, а все внутренние уровни положительны, проверка должна
перейти к высокоугловому хвосту полного тензорного оператора.

\begin{protocol}
\begin{enumerate}
 \item дополнительный профиль: \textbf{$q(r)$};
 \item новые свободные параметры: \textbf{нет};
 \item значение $q(0)$: \textbf{$0.9666849845\ldots$};
 \item полный стационарный фон: \textbf{построен};
 \item натяжение: \textbf{$1.5744526907\ldots$};
 \item вириальный остаток: \textbf{$6.27\cdot10^{-9}$};
 \item полный скрученный спектр на новом фоне: \textbf{не проверен};
 \item полная устойчивость вихря: \textbf{не закрыта};
 \item следующий гейт: \textbf{стационарный скрученный спектр}.
\end{enumerate}
\end{protocol}

Машинный сертификат сохранён в
\path{s2t_v6_bosonic_defect_full_tensor_stationary_background_gate_results.json}.

\begin{thebibliography}{9}\raggedright
\bibitem{stationary-background-alonso}
J.~M.~Alonso-Izquierdo, W.~Garc\'ia Fuertes, J.~Mateos Guilarte,
\emph{Dissecting zero modes and bound states on BPS vortices in
Ginzburg--Landau superconductors}, arXiv:1602.09084.

\bibitem{stationary-background-a4-vortex}
\emph{Non-Abelian $A_4$ vortices in $SO(3)$ gauge theory and
non-invertible symmetries}, arXiv:2607.12366.
\end{thebibliography}